\chapter{ALJABAR LINEAR DAN TEOREMA SPEKTRAL}

Dalam mata kuliah analisis tingkat awal, penekanan diberikan pada perilaku fungsi-fungsi individual (bernilai real atau kompleks). Dalam mata kuliah \emph{analisis fungsional},
fokus beralih ke \emph{ruang-ruang} fungsi tersebut dan kelas-kelas pemetaan tertentu di antara ruang-ruang itu. Ternyata uraian yang ringkas,
mungkin terlampau menyederhanakan, tetapi jelas tidak menyesatkan, tentang materi tersebut ialah \emph{aljabar linear berdimensi tak hingga}. Dari sudut pandang
seorang analis, pencapaian terbesar aljabar linear sebagian besar berupa teorema-teorema tentang operator pada ruang vektor berdimensi hingga
(terutama $\R^n$ dan $\C^n$). Akan tetapi, ruang vektor pemetaan bernilai skalar yang didefinisikan pada berbagai domain biasanya berdimensi tak
hingga. Apakah proses memperumum hasil-hasil mengagumkan dari aljabar linear berdimensi hingga ke ranah berdimensi tak hingga dipandang sebagai
rangkaian komplikasi abstrak yang sangat tidak menyenangkan atau sebagai sumber kaya wawasan menakjubkan, pertanyaan baru yang menarik, dan
petunjuk yang menggugah menuju penerapan di bidang lain sepenuhnya bergantung pada orientasi seseorang terhadap matematika secara umum.

Berdasarkan uraian di atas, akan bermanfaat untuk memulai perjalanan kita ke wilayah baru ini dengan tinjauan singkat atas beberapa keberhasilan
klasik aljabar linear. Tidak jarang mahasiswa yang berhasil menyelesaikan mata kuliah pengantar aljabar linear pada akhirnya hanya memiliki
gagasan yang paling samar tentang pokok bahasan mata kuliah itu. Sebagian kesalahan mungkin terletak pada banyak buku dasar yang mencampuradukkan
secara tidak semestinya dua pokok bahasan yang sangat berbeda, meskipun berkaitan erat: di satu pihak, kajian \emph{ruang vektor} dan pemetaan linear
di antaranya, dan di pihak lain, \emph{ruang hasil kali dalam} beserta pemetaan linearnya. Ruang vektor tidak memiliki gagasan geometri/topologi
tentang \emph{jarak}, \emph{panjang}, \emph{ketegaklurusan}, himpunan \emph{terbuka}, atau \emph{sudut} antara vektor; yang ada hanyalah penjumlahan dan perkalian skalar.
Ruang hasil kali dalam adalah ruang vektor yang dilengkapi struktur-struktur tambahan tersebut. Mari kita tinjau keduanya secara terpisah.


\section{Ruang Vektor dan Dekomposisi Operator yang Dapat Didiagonalkan}
\begin{conv} Dalam catatan ini semua ruang vektor
 \index{konvensi!ruang vektor kompleks atau real}%
 \index{C@$\C$!medan bilangan kompleks}%
 \index{R@$\R$!medan bilangan real}%
akan diasumsikan sebagai ruang vektor atas medan~$\C$ bilangan kompleks (dalam hal ini disebut \emph{ruang vektor kompleks}) atau atas
medan~$\R$ bilangan real (dalam hal ini disebut \emph{ruang vektor real}). Tidak ada medan lain yang akan muncul. Ketika $\K$ muncul, medan itu dapat dipandang
 \index{K@$\K$ (medan bilangan real atau kompleks)}%
sebagai $\C$ ataupun~$\R$. Suatu
 \index{skalar}%
\df{skalar} adalah anggota $\K$; yaitu, suatu \emph{bilangan kompleks} atau suatu \emph{bilangan real}.
\end{conv}

\begin{defn}\label{00001} Tripel $(V,+,M)$ adalah
 \index{vektor!ruang}%
 \index{ruang!vektor}%
\df{ruang vektor} atas $\K$ jika $(V,+)$ adalah grup Abelian dan $M\colon \K \sto \Hom(V)$
adalah homomorfisme gelanggang beridentitas (dengan $\Hom(V)$ gelanggang homomorfisme grup
pada~$V$).
\end{defn}

Untuk memeriksa arti istilah-istilah dalam definisi sebelumnya, lihat tiga bagian pertama dari bab pertama catatan
aljabar linear saya~\cite{Erdman:2010}.

\begin{exer} Definisi \emph{ruang vektor} yang ditemukan dalam banyak buku dasar kurang lebih
sebagai berikut: suatu \emph{ruang vektor atas} $\K$ adalah himpunan $V$ beserta operasi
penjumlahan dan perkalian skalar yang memenuhi aksioma-aksioma berikut:
 \begin{enumerate}
  \item[(1)] jika $x$, $y \in V$, maka $x + y \in V$;
  \item[(2)] $(x + y) + z = x + (y + z)$ untuk setiap
 \index{asosiatif}%
$x$, $y$, $z \in V$ (asosiativitas);
  \item[(3)] terdapat $\vc 0 \in V$ sehingga $x + \vc 0 = x$ untuk setiap
 \index{identitas!aditif}%
$x \in V$ (keberadaan identitas aditif);
  \item[(4)] untuk setiap $x\in V$ terdapat $-x \in V$ sehingga
 \index{aditif!invers}%
 \index{invers!aditif}%
$x + (-x) = \vc 0$ (keberadaan invers aditif);
  \item[(5)] $x + y = y + x$ untuk setiap $x$, $y \in V$
 \index{komutatif}%
(komutativitas);
  \item[(6)] jika $\alpha \in \K$ dan $x \in V$, maka $\alpha x \in V$;
  \item[(7)] $\alpha (x + y) = \alpha x + \alpha y$ untuk setiap $\alpha \in \K$
dan setiap $x$, $y \in V$;
  \item[(8)] $(\alpha + \beta)x = \alpha x + \beta x$ untuk setiap $\alpha$,
$\beta \in \K$ dan setiap $x \in V$;
  \item[(9)] $(\alpha\beta)x = \alpha(\beta x)$ untuk setiap $\alpha$,
$\beta \in \K$ dan setiap $x \in V$; dan
  \item[(10)] $1\,x = x$ untuk setiap $x \in V$.
 \end{enumerate}
Buktikan bahwa definisi ini ekuivalen dengan definisi yang diberikan di atas dalam~\ref{00001}.
\end{exer}

\begin{defn} Subhimpunan $M$ dari ruang vektor $V$ adalah
 \index{subruang!vektor}%
 \index{vektor!subruang}%
\df{subruang vektor} dari $V$ jika merupakan ruang vektor terhadap operasi yang diwarisinya dari~$V$.
\end{defn}

\begin{prop} Subhimpunan tak kosong $M$ dari ruang vektor $V$ adalah subruang vektor dari $V$ jika dan
hanya jika tertutup terhadap penjumlahan dan perkalian skalar. (Artinya: jika $\vc x$ dan
$\vc y$ termasuk dalam $M$, maka $\vc x + \vc y$ juga demikian; dan jika $\vc x$ termasuk dalam $M$ dan $\alpha
\in \K$, maka $\alpha \vc x$ termasuk dalam~$M$.)
\end{prop}

\begin{defn} Vektor $y$ adalah
 \index{linear!kombinasi}%
 \index{kombinasi!linear}%
\df{kombinasi linear} dari vektor-vektor $x_1$, \dots, $x_n$ jika terdapat skalar $\alpha_1$,
\dots, $\alpha_n$ sehingga $y = \sum_{k=1}^n \alpha_k x_k$. Kombinasi linear $\sum_{k=1}^n \alpha_k x_k$ disebut
 \index{linear!kombinasi!trivial}%
 \index{trivial!kombinasi linear}%
\df{trivial} jika semua koefisien $\alpha_1$, \dots, $\alpha_n$ bernilai nol. Jika sekurang-kurangnya
satu $\alpha_k$ tidak sama dengan nol, kombinasi linear itu disebut \df{nontrivial}.
\end{defn}

\begin{defn} Jika $A$ adalah subhimpunan tak kosong dari ruang vektor~$V$, maka $\spn A$, yaitu
 \index{rentang}%
\df{rentang} dari $A$, adalah himpunan semua kombinasi linear unsur-unsur~$A$. Himpunan ini juga sering disebut
 \index{rentang!linear}%
 \index{linear!rentang}%
\df{rentang linear} dari~$A$.
\end{defn}

\begin{notn}\label{bases111} Misalkan $A$ adalah himpunan yang merentang ruang vektor~$V$. Untuk setiap
$x \in V$ terdapat himpunan hingga $S$ yang terdiri atas vektor-vektor dalam $A$, dan untuk setiap unsur $e$ dalam $A$
terdapat skalar $x_e$ sehingga $x = \sum_{e \in S} x_e e$. Jika kita sepakat menetapkan $x_e = 0$
apabila $e \in A \setminus S$, kita dapat pula menulis $x = \sum_{e \in A} x_e e$. Walaupun
notasi ini mungkin memberi kesan bahwa kita menjumlahkan atas himpunan sembarang yang mungkin tak terhitung,
kenyataannya semua kecuali sejumlah hingga suku bernilai nol, sehingga tidak timbul masalah
``kekonvergenan''. Perlakukan $\sum_{e \in A} x_e e$ sebagai jumlah hingga.
Asosiativitas dan komutativitas penjumlahan dalam $V$ membuat ekspresi tersebut tidak ambigu.
\end{notn}

\begin{defn}\label{defn_lin_dep} Subhimpunan $A$ dari suatu ruang vektor disebut
 \index{linear!kebergantungan}%
 \index{bergantung}
\df{bergantung linear} jika vektor nol $\vc 0$ dapat ditulis sebagai kombinasi linear nontrivial
dari unsur-unsur~$A$; yaitu, jika terdapat vektor $x_1, \dots, x_n \in A$
dan skalar $\alpha_1, \dots, \alpha_n$, yang \textbf{tidak semuanya nol,} sehingga $\sum_{k=1}^n
\alpha_kx_k = \vc 0$. Subhimpunan suatu ruang vektor disebut
 \index{linear!kebebasan}%
 \index{bebas}
\df{bebas linear} jika tidak bergantung linear.
\end{defn}

Secara teknis, suatu \emph{himpunan} vektorlah yang bergantung atau bebas linear.
Meskipun demikian, istilah-istilah ini sering digunakan seolah-olah merupakan sifat vektornya
sendiri. Sebagai contoh, jika $S = \{x_1, \dots, x_n\}$ adalah himpunan hingga vektor dalam suatu
ruang vektor, pernyataan ``himpunan $S$ bebas linear'' dan
``vektor-vektor $x_1$, \dots, $x_n$ bebas linear'' dapat digunakan secara bergantian.

\begin{defn}\label{defn_Hamel_basis} Himpunan $B$ yang terdiri atas vektor-vektor dalam ruang vektor $V$ adalah
 \index{basis!Hamel}%
 \index{Hamel!basis}%
\df{basis Hamel} untuk $V$ jika bebas linear dan merentang~$V$.
\end{defn}

\begin{conv}\label{conv_1_1} Ruang vektor yang hanya memuat satu vektor, yaitu vektor nol, adalah ruang vektor
 \index{konvensi!basis untuk ruang vektor nol}%
 \index{vektor!ruang!trivial (atau nol)}%
 \index{trivial!ruang vektor}%
\df{trivial} (atau
 \index{nol!ruang vektor}%
\df{nol}). Terkait ruang ini terdapat sebuah pertanyaan teknis, meskipun tidak terlalu menarik. Apakah ruang ini memiliki basis Hamel?
Dari definisi \emph{kebergantungan linear}~\ref{defn_lin_dep}, jelas bahwa himpunan kosong $\emptyset$ bebas
linear. Himpunan kosong tentu merupakan subhimpunan bebas linear maksimal dari ruang nol (suatu kondisi yang, untuk ruang nontrivial,
ekuivalen dengan menjadi himpunan perentang yang bebas linear). Akan tetapi, sulit untuk berpendapat bahwa himpunan kosong merentang sesuatu.
Jadi, semata-mata sebagai konvensi, kita akan mengatakan bahwa jawaban atas pertanyaan itu adalah \emph{ya}. Ruang vektor nol memiliki basis
Hamel, dan basis itu adalah himpunan kosong!
\end{conv}

Dalam buku pengantar aljabar linear, himpunan perentang yang bebas linear hanya disebut \emph{basis}. Di sini digunakan istilah \emph{basis Hamel}
untuk membedakannya dari \emph{basis ortonormal} dan \emph{basis Schauder}, istilah-istilah yang mengacu pada konsep yang lebih
penting dalam analisis fungsional.

\begin{prop} Misalkan $B$ adalah basis Hamel untuk ruang vektor nontrivial~$V$. Maka setiap unsur dalam $V$ dapat ditulis secara tunggal sebagai kombinasi
linear anggota-anggota~$B$.
\end{prop}

\begin{proof}[\emph{Petunjuk pembuktian}] Eksistensinya jelas. Sederhanakan pembuktian ketunggalan dengan menggunakan konvensi notasi
yang disarankan dalam~\ref{bases111}.   \ns
\end{proof}

\begin{prop}\label{bases025} Misalkan $A$ adalah subhimpunan bebas linear dari ruang
vektor~$V$. Maka terdapat basis Hamel untuk~$V$ yang memuat~$A$.
\end{prop}

\begin{proof}[\emph{Petunjuk pembuktian}] Tunjukkan terlebih dahulu bahwa, untuk ruang nontrivial, subhimpunan bebas linear dari $V$ adalah basis untuk~$V$
jika dan hanya jika merupakan subhimpunan bebas linear maksimal. Kemudian urutkan keluarga subhimpunan bebas linear dari
$V$ yang memuat~$A$ berdasarkan inklusi dan terapkan \emph{lema Zorn}. (Jika Anda belum mengenal \emph{lema Zorn},
lihat Bagian 1.7 dalam catatan aljabar linear saya~\cite{Erdman:2010}.)  \ns
\end{proof}

\begin{cor} Setiap ruang vektor memiliki basis Hamel.
\end{cor}

\begin{defn}\label{000082} Misalkan $M$ dan $N$ adalah subruang dari ruang vektor $V$. Jika
$M \cap N = \{\vc 0\}$ dan $M + N = V$, maka $V$ adalah
 \index{internal!jumlah langsung!dalam ruang vektor}%
 \index{langsung!jumlah!internal}%
 \index{langsung!jumlah!ruang vektor}%
 \index{<binopdirectsum@$\oplus$ (jumlah langsung ruang vektor)}%
\df{jumlah langsung} (\df{internal}) dari $M$ dan $N$. Dalam hal ini kita menulis
   \[ V = M \oplus N\,. \]
Kita mengatakan bahwa $M$ dan $N$ adalah
 \index{komplemen!subruang vektor}%
 \index{vektor!ruang!komplemen}%
\df{subruang vektor komplementer} dan masing-masing merupakan \df{komplemen} (ruang vektor) dari
yang lain. 
 \index{kodimensi}%
\df{Kodimensi} subruang $M$ adalah dimensi komplemennya~$N$.
\end{defn}

\begin{prop}\label{0000823} Misalkan $V$ adalah ruang vektor dan andaikan $V = M \oplus N$. Maka untuk
setiap $v \in V$ terdapat vektor-vektor tunggal $m \in M$ dan $n \in N$ sehingga $v = m + n$.
\end{prop}

\begin{exam} Misalkan $\fml C = \fml C[-1,1]$ adalah ruang vektor semua fungsi kontinu bernilai real
pada interval $[-1,1]$. Fungsi $f$ dalam $\fml C$ disebut \df{genap} jika $f(-x) =
f(x)$ untuk semua $x \in [-1,1]$; fungsi itu disebut \df{ganjil} jika $f(-x) = -f(x)$ untuk semua $x \in [-1,1]$.
Misalkan $\fml C_o = \{f \in \fml C\colon  f \text{ ganjil }\}$ dan $\fml C_e = \{f \in \fml
C\colon  f \text{ genap }\}$. Maka $\fml C = \fml C_o \oplus \fml C_e$.
\end{exam}

\begin{prop}\label{0000824} Jika $M$ adalah subruang dari ruang vektor $V$, maka terdapat subruang $N$ dari
$V$ sehingga $V = M \oplus N$.
\end{prop}

\begin{lem}\label{bases030} Misalkan $V$ adalah ruang vektor dengan basis hingga tak kosong $\{e^1,
\dots,e^n\}$ dan misalkan $v = \sum_{k=1}^n \alpha_k e^k$ adalah vektor dalam~$V$. Jika $p \in \N_n$
dan $\alpha_p \ne 0$, maka
\[
\{e^1, \dots,e^{p-1},v,e^{p+1}, \dots, e^n\}
\]
adalah basis untuk~$V$.
\end{lem}

\begin{prop}\label{bases031} Jika suatu basis untuk ruang vektor $V$ memuat $n$ unsur,
maka setiap subhimpunan bebas linear dari $V$ yang memiliki $n$ unsur juga merupakan basis.
\end{prop}

\begin{proof}[\emph{Petunjuk pembuktian}] Andaikan $\{e^1,\dots,e^n\}$ adalah basis untuk~$V$ dan
$\{v^1, \dots,v^n\}$ bebas linear dalam~$V$. Mulailah dengan menggunakan lema~\ref{bases030}
untuk menunjukkan bahwa (setelah mungkin menomori ulang $e^k$) himpunan $\{v^1, e^2, \dots, e^n\}$
adalah basis untuk~$V$.  \ns
\end{proof}

\begin{cor}\label{bases032} Jika ruang vektor $V$ memiliki basis hingga $B$, maka setiap basis
untuk $V$ berhingga dan memuat jumlah unsur yang sama dengan~$B$.
\end{cor}

\begin{defn} Suatu ruang vektor disebut
 \index{berdimensi hingga}%
\df{berdimensi hingga} jika memiliki basis hingga, dan
 \index{dimensi!ruang vektor}%
 \index{vektor!ruang!dimensi}%
 \index{dim@$\dim V$ (dimensi ruang vektor $V$)}%
\df{dimensi} ruang tersebut adalah banyaknya unsur dalam basis ini (dan dengan demikian dalam basis mana pun) untuk
ruang itu. Dimensi ruang vektor berdimensi hingga $V$ dinotasikan dengan $\dim V$. Berdasarkan
konvensi yang dibuat di atas~\ref{conv_1_1}, ruang vektor nol berdimensi nol. Jika
ruang tersebut tidak memiliki basis hingga, ruang itu disebut
 \index{berdimensi tak hingga}%
\df{berdimensi tak hingga}.
\end{defn}

\begin{defn}\label{defn_op_vsp} Fungsi $T\colon V \sto W$ antara ruang-ruang vektor (atas medan yang sama) disebut
 \index{linear!pemetaan}%
\df{linear} jika $T(u + v) = Tu + Tv$ untuk semua $u$, $v \in V$ dan $T(\alpha v) = \alpha Tv$
untuk semua $\alpha \in \K$ dan $v \in V$. Fungsi linear sering disebut
\emph{transformasi linear} atau \emph{pemetaan linear}. Jika $V = W$, kita mengatakan bahwa pemetaan linear $T$ adalah
 \index{operator!pada ruang vektor}%
\df{operator} pada~$V$. Bergantung pada konteks, operator identitas $x \mapsto x$ pada $V$ dinotasikan
 \index{identitas@$\id V$ atau $I_V$ atau $I$ (operator identitas)}%
dengan $\id V$ atau $I_V$ atau cukup~$I$. Ingat bahwa jika $T\colon V \sto W$ adalah pemetaan linear, maka
 \index{kernel!pemetaan linear}%
 \index{<ker@$\ker T$ (kernel dari~$T$)}%
\df{kernel} dari $T$, yang dinotasikan dengan $\ker T$, adalah $T^{\gets}(\{0\}) := \{x \in V\colon Tx = \vc
0\}$. Selain itu,
 \index{jangkauan!pemetaan linear}%
 \index{<ran@$\ran T$ (jangkauan dari~$T$)}%
\df{jangkauan} dari $T$, yang dinotasikan dengan $\ran T$, adalah $T^{\sto}(V) := \{Tx\colon x \in V\}$.
\end{defn}

\begin{prop} Pemetaan linear $T\colon V \sto W$ bersifat injektif jika dan hanya jika kernelnya nol.
\end{prop}

\begin{prop} Misalkan $T\colon V \sto W$ adalah pemetaan linear injektif antara ruang-ruang vektor. Jika $A$ adalah subhimpunan bebas linear
dari~$V$, maka $T^\sto(A)$ adalah subhimpunan bebas linear dari~$W$.
\end{prop}

\begin{prop} Misalkan $T\colon V \sto W$ adalah pemetaan linear antara ruang-ruang vektor dan $A \subseteq V$.
Maka $T^\sto(\spn A) = \spn T^\sto(A)$.
\end{prop}

\begin{prop} Misalkan $T\colon V \sto W$ adalah pemetaan linear injektif antara ruang-ruang vektor. Jika $B$ adalah basis
untuk subruang $U$ dari~$V$, maka $T^\sto(B)$ adalah basis untuk~$T^\sto(U)$.
\end{prop}

\begin{prop} Dua ruang vektor berdimensi hingga isomorfik jika dan hanya jika keduanya memiliki dimensi yang sama.
\end{prop}

\begin{defn} Misalkan $T$ suatu pemetaan linear antara ruang-ruang vektor. \df{Peringkat} $T$ adalah dimensi jangkauan~$T$, sedangkan
 \index{peringkat}
 \index{nulitas}%
\df{nulitas} $T$ adalah dimensi kernel~$T$.
\end{defn}

\begin{prop} Misalkan $T\colon V \sto W$ adalah pemetaan linear antara ruang-ruang vektor. Jika $V$ berdimensi hingga, maka
  \[ \operatorname{peringkat} T + \operatorname{nulitas} T = \dim V. \]
\end{prop}

\begin{defn} Pemetaan linear $T\colon V \sto W$ antara ruang-ruang vektor disebut
 \index{invertibel!pemetaan linear}%
\df{invertibel} (atau suatu
 \index{isomorfisme!ruang vektor}%
\df{isomorfisme}) jika terdapat pemetaan linear $T^{-1}\colon W \sto V$ sehingga $T^{-1}T
= \id V$ dan $TT^{-1} = \id W$. Dua ruang vektor $V$ dan $W$ disebut
 \index{isomorfik!ruang vektor}%
\df{isomorfik} jika terdapat isomorfisme dari $V$ ke~$W$.
\end{defn}

\begin{prop}\label{inv_lin_maps001} Pemetaan linear $T\colon V \sto W$ antara ruang-ruang vektor
invertibel jika dan hanya jika memiliki invers kiri sekaligus invers kanan.
\end{prop}

\begin{prop}\label{inv_lin_maps004} Pemetaan linear antara ruang-ruang vektor invertibel jika
dan hanya jika bijektif.
\end{prop}

\begin{prop} Misalkan $T$ operator pada ruang vektor berdimensi hingga. Ketiga pernyataan berikut ekuivalen:
 \begin{enumerate}
    \item $T$ adalah isomorfisme;
    \item $T$ injektif; dan
    \item $T$ surjektif.
 \end{enumerate}
\end{prop}

 \begin{defn}\label{defn_similar_vs_ops} Dua operator $R$ dan $T$ pada ruang vektor $V$ disebut
 \index{serupa!operator}%
\df{serupa} jika terdapat operator invertibel $S$ pada $V$ sehingga $R = STS^{-1}$.
\end{defn}

\begin{prop}\label{sim004thm} Jika $V$ adalah ruang vektor, maka keserupaan merupakan relasi ekuivalensi
pada~$\ofml L(V)$.
\end{prop}

Misalkan $V$ dan $W$ adalah ruang vektor berdimensi hingga dengan basis terurut. Andaikan $V$ berdimensi $n$ dengan basis terurut
$\{e^1,e^2, \dots,e^n\}$ dan $W$ berdimensi $m$. Ingat dari pengantar aljabar linear bahwa jika $T\colon V \sto W$ linear,
maka
 \index{matriks!representasi}%
 \index{representasi!matriks}%
\emph{representasi matriks} $[T]$ (diambil terhadap basis terurut dalam $V$ dan~$W$) adalah matriks $m \times n$ $[T]$
yang kolom ke-$k$ ($1 \le k \le n$) adalah vektor kolom $Te^k$ dalam~$W$. Pokoknya ialah bahwa tindakan $T$ pada
vektor $x$ dapat \emph{direpresentasikan} sebagai perkalian $x$ oleh matriks $[T]$; yaitu,
   \[ Tx = [T]x. \]
Persamaan ini mungkin memerlukan sedikit penafsiran. Ruas kiri adalah fungsi $T$ yang dievaluasi di $x$, dengan
hasil evaluasi dipandang sebagai vektor kolom dalam~$W$; ruas kanan adalah matriks $m \times n$ yang dikalikan
dengan matriks $n \times 1$ (yaitu vektor kolom). Jadi, kesamaan yang dinyatakan adalah kesamaan dua matriks $m \times 1$ (vektor kolom).

\begin{defn} Misalkan $V$ adalah ruang vektor berdimensi hingga dan $B = \{e^1, \dots, e^n\}$ adalah basis untuk~$V$. Operator $T$ pada $V$ disebut
 \index{diagonal!operator}%
 \index{operator!diagonal}%
\df{diagonal} jika terdapat skalar $\alpha_1, \dots, \alpha_n$ sehingga $Te^k =
\alpha_ke^k$ untuk setiap $k \in \N_n$. Secara ekuivalen, $T$ diagonal jika representasi matriksnya
$[T] = [T_{ij}]$ memiliki sifat bahwa $T_{ij} = 0$ apabila $i \ne j$.
\end{defn}

Menanyakan apakah operator tertentu pada suatu ruang vektor berdimensi hingga bersifat diagonal,
secara tegas, tidak bermakna. Sebagaimana didefinisikan, sifat operator berupa kediagonalan jelas
\emph{bukan} konsep ruang vektor. Sifat ini hanya bermakna bagi ruang vektor
\emph{yang basisnya telah ditentukan}. Fakta penting ini, meskipun jelas, tampaknya
luput dari perhatian dalam banyak mata kuliah pengantar aljabar linear, mungkin karena fiksasi yang agak berlebihan
pada $\R^n$ dalam mata kuliah tersebut. Berikut sifat \emph{ruang vektor} yang relevan.

\begin{defn} Operator $T$ pada ruang vektor berdimensi hingga $V$ disebut
 \index{dapat didiagonalkan!operator}%
 \index{operator!dapat didiagonalkan}%
\df{dapat didiagonalkan} jika terdapat basis untuk $V$ yang terhadapnya $T$
bersifat diagonal. Secara ekuivalen, operator pada ruang vektor berdimensi hingga \emph{dengan
basis} dapat didiagonalkan jika serupa dengan operator diagonal.
\end{defn}

\begin{defn} Misalkan $V$ adalah ruang vektor dan andaikan $V = M \oplus N$. Kita mengetahui
dari~\ref{0000823} bahwa untuk setiap $\vc v \in V$ terdapat vektor-vektor tunggal $\vc m \in M$
dan $\vc n \in N$ sehingga $\vc v = \vc m + \vc n$. Definisikan fungsi $\sbsb E{MN}\colon V \sto V$ dengan $\sbsb E{MN}v = n$.
Fungsi $\sbsb E{MN}$ adalah
 \index{proyeksi!dalam ruang vektor}%
 \index{E@$\sbsb E{MN}$ (proyeksi sepanjang $M$ ke $N$)}%
\df{proyeksi} $V$ \df{sepanjang} $M$ \df{ke}~$N$. (Kita sering menulis $E$ untuk $\sbsb E{MN}$. Namun ingat bahwa
$E$ bergantung pada $M$ dan~$N$.)
\end{defn}

\begin{prop} Misalkan $V$ adalah ruang vektor dan andaikan $V = M \oplus N$. Jika $E$ adalah
proyeksi $V$ sepanjang $M$ ke $N$, maka
 \begin{enumerate}[label=\rm{(\alph*)}]
  \item $E$ linear;
 \index{idempoten}%
  \item $E^2 = E$ \quad (yaitu, $E$ \df{idempoten});
  \item $\ran E = N$; dan
  \item $\ker E = M$.
 \end{enumerate}
\end{prop}

\begin{prop} Misalkan $V$ adalah ruang vektor dan andaikan $E\colon V \sto V$ adalah fungsi yang
memenuhi
 \begin{enumerate}[label=\rm{(\alph*)}]
  \item $E$ linear, dan
  \item $E^2 = E$.
\end{enumerate}
Maka
 \[V = \ker E \oplus \ran E\]
dan $E$ adalah proyeksi $V$ sepanjang $\ker E$ ke~$\ran E$.
\end{prop}

Penting untuk dicatat bahwa konsekuensi langsung dari dua proposisi terakhir adalah bahwa
fungsi $T\colon V \sto V$ dari ruang vektor berdimensi hingga ke dirinya sendiri merupakan
proyeksi jika dan hanya jika linear dan idempoten.






\begin{prop} Misalkan $T\colon V \sto W$ adalah transformasi linear antara ruang-ruang vektor. Maka
  \begin{enumerate}[label=\rm{(\alph*)}]
    \item $T$ memiliki invers kiri jika dan hanya jika injektif; dan
    \item $T$ memiliki invers kanan jika dan hanya jika surjektif.
  \end{enumerate}
\end{prop}











\begin{prop} Misalkan $V$ adalah ruang vektor dan andaikan $V = M \oplus N$. Jika $E$ adalah
proyeksi $V$ sepanjang $M$ ke $N$, maka $I - E$ adalah proyeksi $V$ sepanjang $N$
ke~$M$.
\end{prop}

Seperti baru saja kita lihat, jika $E$ adalah proyeksi pada ruang vektor $V$, maka operator
identitas pada $V$ dapat ditulis sebagai jumlah dua proyeksi $E$ dan $I - E$ yang
jangkauannya membentuk dekomposisi jumlah langsung ruang tersebut, $V = \ran E \oplus \ran
(I - E)$. Kita dapat memperumum hal ini ke lebih dari dua proyeksi.

\begin{defn} Andaikan pada ruang vektor $V$ terdapat operator-operator proyeksi $E_1$,
\dots, $E_n$ sehingga
\begin{enumerate}
    \item $I_V = E_1 + E_2 + \dots + E_n$ dan
    \item $E_iE_j = 0$ apabila $i \ne j$.
\end{enumerate}
Kita mengatakan bahwa $I_V = E_1 + E_2 + \dots + E_n$ adalah
 \index{resolusi identitas!dalam ruang vektor}%
\df{resolusi identitas}.
\end{defn}

\begin{prop} Jika $I_V = E_1 + E_2 + \dots + E_n$ adalah resolusi identitas pada ruang
vektor $V$, maka $V = \bigoplus_{k=1}^n \ran E_k$.
\end{prop}

\begin{exam} Misalkan $P$ adalah bidang dalam $\R^3$ yang persamaannya $x - z = 0$ dan $L$ adalah garis
yang persamaannya $y = 0$ dan $x = -z$. Misalkan $E$ adalah proyeksi $\R^3$ sepanjang~$L$
ke $P$ dan $F$ adalah proyeksi $\R^3$ sepanjang~$P$ ke~$L$. Maka
 \[[E] = \begin{bmatrix}
              \frac12  & 0 & \frac12 \\
                  0    & 1 &   0     \\
              \frac12  & 0 & \frac12
         \end{bmatrix} \qquad \text{ dan } \qquad
   [F] = \begin{bmatrix}
              \frac12  & 0 & -\frac12 \\
                  0    & 0 &   0     \\
             -\frac12  & 0 &  \frac12
         \end{bmatrix}\,. \]
\end{exam}

\begin{defn} Skalar $\lambda \in \K$ adalah
 \index{nilai eigen}%
\df{nilai eigen} dari operator $T$ pada ruang vektor $V$ jika $\ker(T - \lambda I_V)$
memuat vektor tak nol. Setiap vektor tersebut adalah
 \index{vektor eigen}%
\df{vektor eigen} dari $T$ yang berkaitan dengan $\lambda$, dan $\ker(T - \lambda I_V)$ adalah
 \index{ruang eigen}
\df{ruang eigen} dari $T$ yang berkaitan dengan~$\lambda$. Himpunan semua nilai eigen
operator $T$ disebut
 \index{spektrum titik}%
 \index{spektrum!titik}%
\df{spektrum titik} dan dinotasikan dengan~$\sigma_p(T)$.
\end{defn}

Jika $M$ adalah matriks $n \times n$, maka $\det(M - \lambda I_n)$ (dengan $I_n$ matriks identitas $n
\times n$) adalah polinomial dalam $\lambda$ berderajat~$n$. Inilah
\df{polinomial karakteristik} dari~$M$. Cara baku untuk menghitung nilai-nilai eigen suatu
operator $T$ pada ruang vektor berdimensi hingga adalah mencari akar-akar polinomial
karakteristik dari representasi matriksnya. Sifat multiplikatif fungsi determinan dengan mudah mengakibatkan
bahwa polinomial karakteristik operator $T$ pada ruang vektor $V$ tidak bergantung pada basis yang dipilih untuk $V$ dan,
karena itu, tidak bergantung pada representasi matriks tertentu dari~$T$ yang digunakan.

\vskip 5 pt

\begin{exam} Nilai-nilai eigen operator pada (ruang vektor real) $\R^3$ yang representasi matriksnya
adalah $ \begin{bmatrix} 0  &  0  &  2 \\
                                     0  &  2  &  0 \\
                                     2  &  0  &  0
                      \end{bmatrix}$ ialah $-2$ dan $+2$, dengan nilai yang terakhir memiliki multiplisitas (aljabar maupun
geometris)~$2$. Ruang eigen yang berkaitan dengan nilai eigen negatif dan positif, berturut-turut, adalah
\[
\spn\{(1,0,-1)\}
\quad\text{dan}\quad
\spn\{(1,0,1), (0,1,0)\}.
\]
\end{exam}

\vskip 5 pt

\begin{prop} Setiap operator pada ruang vektor kompleks berdimensi hingga memiliki nilai eigen.
\end{prop}

\begin{exam} Proposisi sebelumnya tidak berlaku untuk ruang real berdimensi hingga.
\end{exam}

\begin{proof}[\emph{Petunjuk pembuktian}] Pertimbangkan rotasi bidang.  \ns
\end{proof}

Fakta utama yang dinyatakan oleh versi ruang vektor berdimensi hingga dari
\emph{teorema spektral} ialah bahwa setiap operator yang dapat didiagonalkan pada ruang demikian dapat
ditulis sebagai kombinasi linear operator-operator proyeksi, dengan koefisien kombinasi
linear tersebut berupa nilai-nilai eigen operator dan jangkauan proyeksinya berupa
ruang eigen yang bersesuaian. Jadi, jika $T$ adalah operator yang dapat didiagonalkan pada ruang vektor
berdimensi hingga $V$, maka $V$ memiliki basis yang terdiri atas vektor-vektor eigen dari~$T$.

Berikut pernyataan formal teorema tersebut.

\begin{thm}[Teorema Spektral: versi ruang vektor berdimensi hingga] Andaikan $T$ adalah
 \index{spektral!teorema!operator yang dapat didiagonalkan pada ruang vektor}%
operator yang dapat didiagonalkan pada ruang vektor berdimensi hingga~$V$. Misalkan $\lambda_1, \dots, \lambda_n$ adalah
nilai-nilai eigen (berbeda) dari~$T$. Maka terdapat resolusi identitas $I_V =
E_1 + \dots + E_n$, dengan untuk setiap $k$ jangkauan proyeksi $E_k$ merupakan
ruang eigen yang berkaitan dengan~$\lambda_k$, dan selanjutnya
     \[ T = \lambda_1E_1 +\dots + \lambda_nE_n\,. \]
\end{thm}

\begin{proof} Pembuktian teorema ini dapat ditemukan dalam \cite{HoffmanK:1971} pada halaman~212. \ns
\end{proof}

\begin{exam} Misalkan $T$ adalah operator pada (ruang vektor real) $\R^2$ yang representasi matriksnya
adalah $\begin{bmatrix} -7 &   8 \\ -16 &  17 \end{bmatrix}$.

\vskip 3 pt

 \begin{enumerate}
  \item Polinomial karakteristik untuk~$T$ adalah
                     $\sbsb cT(\lambda) = \lambda^2 - 10\lambda + 9$.

\vskip 3 pt

  \item Ruang eigen $M_1$ yang berkaitan dengan nilai eigen $1$ adalah $\spn \{(1,1)\}$.

\vskip 3 pt

  \item Ruang eigen $M_2$ yang berkaitan dengan nilai eigen $9$ adalah $\spn \{(1,2)\}$.

\vskip 3 pt

  \item Kita dapat menulis $T$ sebagai kombinasi linear operator-operator proyeksi. Secara khusus,
   \[ T = 1\cdot E_1 + 9\cdot E_2 \text{ dengan } [E_1] = \begin{bmatrix} 2 & -1 \\ 2 & -1 \end{bmatrix}
\text{ dan } [E_2] = \begin{bmatrix} -1 & 1 \\ -2 & 2 \end{bmatrix}\,. \]

\vskip 3 pt

  \item Perhatikan bahwa jumlah $[E_1]$ dan $[E_2]$ adalah matriks identitas dan hasil kali keduanya
adalah matriks nol.

\vskip 3 pt

  \item Matriks $S = \begin{bmatrix} 1 &  1  \\ 1 &  2 \end{bmatrix}$ mendiagonalkan
$[T]$. Artinya, $S^{-1}\,[T]\,S = \begin{bmatrix} 1 &  0  \\ 0 &  9 \end{bmatrix}$.


\vskip 3 pt


  \item Matriks yang merepresentasikan $\sqrt T$ adalah $\begin{bmatrix} -1 &  2  \\ -4 &  5 \end{bmatrix}$.
 \end{enumerate}
\end{exam}

\begin{exer} Misalkan $T$ adalah operator pada $\R^3$ yang representasi matriksnya
    \[
      \begin{bmatrix}
               2  &  0  &  0 \\
              -1  &  3  &  2 \\
               1  & -1  &  0
      \end{bmatrix}.
    \]
  \begin{enumerate}
    \item Tentukan polinomial karakteristik dan polinomial minimal dari~$T$.
    \item Apa yang dapat disimpulkan dari bentuk polinomial minimal tersebut?
    \item Tentukan matriks yang mendiagonalkan~$T$. Apa bentuk diagonal $T$ yang dihasilkan oleh matriks ini?
    \item Tentukan (representasi matriks dari)~$\sqrt T$.
  \end{enumerate}
\end{exer}

\begin{defn} Operator ruang vektor $T$ disebut
 \index{nilpoten}%
\df{nilpoten} jika $T^n = \vc 0$ untuk suatu $n \in \N$.
\end{defn}

Operator pada ruang vektor berdimensi hingga tidak harus dapat didiagonalkan. Jika tidak,
seberapa dekat operator itu dengan operator yang dapat didiagonalkan? Berikut salah satu jawabannya.

\begin{thm} Misalkan $T$ adalah operator pada ruang vektor berdimensi hingga~$V$. Andaikan polinomial minimal
untuk $T$ terfaktor sepenuhnya menjadi faktor-faktor linear
   \[ m_T(x) = (x - \lambda_1)^{r_1} \dots (x - \lambda_k)^{r_k} \]
dengan $\lambda_1, \dots \lambda_k$ merupakan nilai-nilai eigen (berbeda) dari~$T$. Untuk setiap $j$, misalkan
$W_j = \ker(T - \lambda_jI)^{r_j}$ dan $E_j$ adalah proyeksi $V$ ke $W_j$ sepanjang
$W_1 + \dots + W_{j-1} + W_{j+1} + \dots + W_k$. Maka
   \[ V = W_1 \oplus W_2 \oplus \dots \oplus W_k, \]
setiap $W_j$ invarian terhadap $T$, dan $I_V = E_1 + \dots + E_k$. Selanjutnya, operator
   \[ D = \lambda_1E_1 + \dots + \lambda_kE_k \]
dapat didiagonalkan, operator
   \[ N = T - D \]
nilpoten, dan $N$ berkomutasi dengan~$D$.
\end{thm}

\begin{proof} Lihat~\cite{HoffmanK:1971}, halaman 222--223. \ns
\end{proof}

\begin{defn} Karena dalam teorema sebelumnya $T = D + N$, dengan $D$ dapat didiagonalkan dan $N$ nilpoten,
kita menyebut $D$ sebagai
 \index{dapat didiagonalkan!bagian}%
 \index{bagian!dapat didiagonalkan}%
\df{bagian yang dapat didiagonalkan} dari $T$, dan $N$ sebagai
 \index{nilpoten!bagian}%
 \index{bagian!nilpoten}%
\df{bagian nilpoten} dari~$T$.
\end{defn}

\begin{exer} Misalkan $T$ adalah operator pada $\R^3$ yang representasi matriksnya
 \[ \begin{bmatrix}
        3  &  1  & -1 \\
        2  &  2  & -1 \\
        2  &  2  &  0
  \end{bmatrix}. \]

 \begin{enumerate}
  \item Tentukan polinomial karakteristik dan polinomial minimal dari~$T$.
  \item Tentukan ruang-ruang eigen dari~$T$.
  \item Tentukan bagian yang dapat didiagonalkan $D$ dan bagian nilpoten $N$ dari~$T$.
  \item Tentukan matriks yang mendiagonalkan~$D$. Apa bentuk diagonal
dari $D$ yang dihasilkan oleh matriks ini?
  \item Tunjukkan bahwa $D$ berkomutasi dengan~$N$.
 \end{enumerate}
\end{exer}
\section{Operator Normal pada Ruang Hasil Kali Dalam}
\begin{defn} Misalkan $V$ suatu ruang vektor. Sebuah fungsi $s$ yang mengaitkan setiap pasangan vektor
$x$ dan $y$ dalam $V$ dengan suatu skalar
 \index{s@$s(x,y)$ (hasil kali dalam semu)}%
 \index{<s@$s(x,y)$ (hasil kali dalam semu)}%
$s(x,y)$ adalah
 \index{hasil kali dalam semu}%
 \index{hasil kali dalam semu!ruang}%
 \index{hasil kali!dalam semu}%
 \index{ruang!hasil kali dalam semu}%
\df{hasil kali dalam semu} pada $V$ apabila untuk setiap $x$, $y$, $z \in V$ dan $\alpha \in \K$ keempat syarat berikut dipenuhi:
 \begin{enumerate}
   \item $s(x + y , z) = s(x,z) + s(y,z)$;
   \item $s(\alpha x,y) = \alpha s(x,y)$;
   \item $s(x,y) = \conj{s(y,x)}$; dan
   \item $s(x,x) \ge 0$.
 \end{enumerate}
 Jika, selain itu, fungsi $s$ memenuhi
  \begin{enumerate}
    \item Untuk setiap $x$ tak nol dalam $V$, berlaku $s(x,x) > 0$.
  \end{enumerate}
maka $s$ adalah
 \index{hasil kali dalam}%
 \index{hasil kali dalam!ruang}%
 \index{hasil kali!dalam}%
 \index{ruang!hasil kali dalam}%
\df{hasil kali dalam} pada~$V$. Hasil kali dalam dua vektor $x$ dan $y$ biasanya akan kita tuliskan sebagai
 \index{<bracpnt@$\langle x,y \rangle$ (hasil kali dalam)}%
$\langle x,y \rangle$, bukan $s(x,y)$.

\noindent Garis atas pada (c) menyatakan konjugasi kompleks (sehingga berlebihan dalam kasus ruang vektor real).
Syarat (a) dan (b) memperlihatkan bahwa hasil kali dalam semu bersifat linear dalam variabel pertamanya. Syarat (a) dan (b) dalam
proposisi~\ref{00015001} menyatakan bahwa hasil kali dalam kompleks bersifat
 \index{konjugat!linear}%
 \index{linear!konjugat}%
\df{linear konjugat} dalam variabel keduanya. Jika suatu fungsi bernilai skalar dengan dua variabel pada ruang hasil kali dalam semu kompleks
bersifat linear dalam satu variabel dan linear konjugat dalam variabel lainnya, fungsi itu sering disebut
 \index{seskuilinear}%
 \index{linear!seskuilinear}%
\df{seskuilinear}. (Awalan ``sesqui-'' berarti ``satu setengah''.) Istilah \emph{seskuilinear} juga akan kita gunakan untuk
fungsi bilinear pada ruang hasil kali dalam semu real. Secara bersama-sama, syarat (a)--(d) menyatakan bahwa hasil kali dalam semu itu merupakan
\emph{bentuk seskuilinear semidefinit positif dan simetris konjugat}; bersama dengan syarat ketat tambahan di atas, bentuk itu definit positif.
\end{defn}

\begin{notn}\label{norm_notn} Jika $V$ adalah ruang vektor yang telah dilengkapi dengan
hasil kali dalam semu dan $x \in V$, kita memperkenalkan singkatan
  \[ \norm x := \sqrt{s(x,x)} \]
yang dibaca \emph{norma $x$} atau \emph{panjang~$x$}. (Istilah yang agak
optimistis ini dibenarkan, setidaknya ketika $s$ merupakan hasil kali dalam, oleh proposisi~\ref{00015025} di bawah.)
\end{notn}

\begin{prop}\label{00015001} Jika $x$, $y$, dan $z$ merupakan vektor dalam suatu ruang dengan hasil kali dalam semu $s$ dan $\alpha \in \K$, maka
 \begin{enumerate}[label=\rm{(\alph*)}]
   \item $s(x,y + z)  = s(x,y) + s(x,z)$,
   \item $s(x,\alpha y) = \conj\alpha s(x,y)$, dan
   \item jika $s$ merupakan hasil kali dalam, maka $s(x,x) = 0$ jika dan hanya jika $x = \vc 0$.
 \end{enumerate}
\end{prop}

\begin{exam}\label{000150012} Untuk vektor $x = (x_1,x_2,\dots,x_n)$ dan $y = (y_1,y_2,\dots,y_n)$
yang termasuk dalam $\K^n$, definisikan
   \[ \langle x,y \rangle = \sum_{k=1}^n x_k \conj{y_k}\,. \]
Maka
 \index{K@$\K^n$!sebagai ruang hasil kali dalam}%
 \index{hasil kali dalam!ruang!$\K^n$ sebagai}%
$\K^n$ merupakan ruang hasil kali dalam.
\end{exam}

\begin{exam}\label{000150013} Misalkan $l_2 = l_2(\N)$ adalah himpunan semua barisan bilangan
kompleks yang terjumlah secara kuadrat. (Suatu barisan $x = (x_k)_{k=1}^\infty$ disebut
 \index{terjumlah secara kuadrat}%
 \index{dapat dijumlahkan!kuadrat}%
 \index{l@$l_2 = l_2(\N)$!barisan terjumlah secara kuadrat}%
 \index{l@$l_2 = l_2(\N)$!sebagai ruang hasil kali dalam}%
 \index{hasil kali dalam!ruang!$l_2$ sebagai}%
\df{terjumlah secara kuadrat} jika $\sum_{k=1}^\infty \abs{x_k}^2 < \infty$.) (Operasi ruang vektornya didefinisikan
titik demi titik.) Untuk vektor $x = (x_1,x_2,\dots)$ dan $y = (y_1,y_2,\dots)$ dalam $l_2$, definisikan
   \[ \langle x,y \rangle = \sum_{k=1}^\infty x_k \conj{y_k}\,. \]
Maka $l_2$ merupakan ruang hasil kali dalam. (Antara lain, harus ditunjukkan bahwa
deret dalam definisi di atas benar-benar konvergen.) Ruang serupa adalah $l_2(\Z)$, yakni himpunan semua
 \index{l@$l_2(\Z)$!barisan dua sisi terjumlah secara kuadrat}%
 \index{l@$l_2(\Z)$!sebagai ruang hasil kali dalam}%
 \index{hasil kali dalam!ruang!$l_2(\Z)$ sebagai}%
barisan dua sisi yang terjumlah secara kuadrat $x = (\dots,x_{-2},x_{-1},x_0,x_1,x_2,\dots)$ dengan hasil kali dalam yang sudah jelas.
\end{exam}

\begin{exam}\label{000150014} Untuk $a < b$, misalkan $\fml C([a,b])$ adalah keluarga semua
 \index{C@$\fml C([a,b])$!sebagai ruang hasil kali dalam}%
 \index{hasil kali dalam!ruang!$\fml C([a,b])$ sebagai}%
fungsi kontinu bernilai kompleks pada interval $[a,b]$. Untuk setiap $f$, $g \in \fml C([a,b])$, definisikan
   \[ \langle f,g \rangle = \int_a^b f(x) \conj{g(x)}\,dx. \]
Maka $\fml C([a,b])$ merupakan ruang hasil kali dalam.
\end{exam}

Berikut adalah fakta tunggal yang paling berguna mengenai hasil kali dalam semu. Fakta ini merupakan suatu ketaksamaan yang secara beragam
dinisbatkan kepada Cauchy, Schwarz, Bunyakowsky, atau gabungan nama ketiganya.

\begin{thm}[Ketaksamaan Schwarz]\label{00015002} Misalkan $s$ suatu hasil kali dalam semu yang didefinisikan pada ruang vektor~$V$.
 \index{ketaksamaan Schwarz}%
 \index{ketaksamaan!Schwarz}%
Maka ketaksamaan
   \[ \abs{s(x,y)} \le \norm x \,\norm y. \]
berlaku untuk semua vektor $x$ dan~$y$ dalam~$V$.
\end{thm}

\begin{proof}[\emph{Petunjuk pembuktian}] Tetapkan $x$, $y \in V$. Untuk setiap $\alpha \in \K$, kita mengetahui bahwa
  \begin{equation}\label{Schwarz_ineq}
      0 \le s(x - \alpha y, x - \alpha y) .
  \end{equation}
Uraikan ruas kanan~\eqref{Schwarz_ineq} menjadi empat suku. Jika $s(y,x)=0$, kesimpulannya langsung. Jika tidak,
tuliskan $s(y,x)$ dalam bentuk polar: $s(y,x) = r e^{i\theta}$, dengan $r > 0$ dan $\theta \in \R$.
Kemudian, dalam ketaksamaan yang diperoleh, tinjau $\alpha$ yang berbentuk $e^{-i\theta}t$ dengan
$t \in \R$. Perhatikan bahwa kini ruas kanan~\eqref{Schwarz_ineq} merupakan polinom kuadrat dalam~$t$. Apa yang dapat Anda katakan tentang
diskriminannya? \ns
\end{proof}

\begin{prop} Misalkan $V$ suatu ruang vektor dengan hasil kali dalam semu $s$ dan $z \in V$. Maka $s(z,z) = 0$
jika dan hanya jika $s(z,y) = 0$ untuk semua $y \in V$.
\end{prop}

\begin{prop} Misalkan $V$ suatu ruang vektor yang telah dilengkapi dengan hasil kali dalam semu $s$. Maka himpunan
$L := \{z \in V\colon s(z,z) = 0\}$ merupakan subruang vektor dari $V$, dan ruang vektor hasil bagi $V/L$ dapat
dijadikan ruang hasil kali dalam dengan mendefinisikan
    \[ \langle\, [x],[y]\, \rangle := s(x,y) \]
untuk semua $[x]$, $[y] \in V/L$.
\end{prop}

Proposisi berikut menunjukkan bahwa hasil kali dalam bersifat kontinu dalam variabel pertamanya. Simetri konjugat kemudian menjamin
kekontinuan dalam variabel keduanya.

\begin{prop} Jika $(x_n)$ merupakan barisan dalam ruang hasil kali dalam $V$ yang konvergen ke suatu
vektor $a \in V$, maka $\langle x_n,y \rangle \sto \langle a,y \rangle$ untuk setiap $y \in V$.
\end{prop}

\begin{exam} Jika $(a_k)$ merupakan barisan bilangan real yang terjumlah secara kuadrat, maka deret $\sum_{k=1}^\infty k^{-1}a_k$
konvergen mutlak.
\end{exam}

\begin{exer} Misalkan $0 < a < b$.
 \begin{enumerate}
  \item Jika $f$ dan $g$ merupakan fungsi kontinu bernilai real pada
interval $[a,b]$, maka
   \[\left(\int_a^b f(x)g(x)\,dx\right)^2
             \le \int_a^b\bigl(f(x)\bigr)^2\,dx \cdot
                       \int_a^b\bigl(g(x)\bigr)^2\,dx.\]

\vskip 3 pt

  \item Gunakan bagian (a) untuk mencari bilangan $M$ dan $N$ (yang bergantung pada $a$
dan~$b$) sedemikian sehingga
   \[M(b-a) \le \ln\left(\frac ba\right) \le N(b-a).\]

\vskip 3 pt

  \item Gunakan bagian (b) untuk menunjukkan bahwa $2.5 \le e \le 3$.
 \end{enumerate}
\end{exer}

\begin{proof}[\emph{Petunjuk pembuktian}] Untuk (b), cobalah $f(x) = \sqrt x$ dan $g(x) = \dfrac1{\sqrt x}$. Kemudian cobalah
$f(x) = \dfrac1x$ dan $g(x) = 1$. Untuk (c), cobalah $a = 1$ dan $b = 3$. Kemudian cobalah $a = 2$ dan $b = 5$. \ns
\end{proof}

\begin{defn} Misalkan $V$ suatu ruang vektor. Sebuah fungsi
 \index{<unopn@$\norm x$ (norma suatu vektor~$x$)}%
$\norm{\hphantom{x}} \colon V \sto \R \colon x \mapsto \norm{x}$ adalah
 \index{norma}%
\df{norma} pada $V$ jika
 \begin{enumerate}
  \item $\norm{x + y} \le \norm{x} + \norm{y}$ \qquad  untuk semua $x$, $y \in V$;
  \item $\norm{\alpha x} = \abs{\alpha}\,\norm{x}$\qquad untuk semua $x \in V$ dan
$\alpha \in \K$; dan
  \item jika $\norm{x} = 0$, maka $x = \vc 0$.
 \end{enumerate}
Ungkapan $\norm{x}$ dapat dibaca sebagai ``\emph{norma} $x$'' atau ``\emph{panjang}
 \index{panjang}%
$x$''. Jika fungsi $\norm{\hphantom{x}}$ memenuhi (a) dan (b)
di atas (tetapi mungkin tidak memenuhi~(c)), fungsi itu merupakan
 \index{seminorma}%
\df{seminorma} pada~$V$.

Ruang vektor yang telah dilengkapi dengan suatu norma disebut
 \index{bernorma!ruang linear}%
 \index{bernorma!ruang vektor|seeonly{ruang linear bernorma}}%
 \index{ruang!linear bernorma}%
\df{ruang linear bernorma} (atau
 \index{vektor!ruang!bernorma}%
\df{ruang vektor bernorma}). Vektor dalam ruang linear bernorma yang normanya $1$ disebut
 \index{satuan!vektor}%
 \index{vektor!satuan}%
\df{vektor satuan}.
\end{defn}

\begin{exer} Tunjukkan mengapa langsung jelas dari definisi bahwa $\norm{\vc 0} = 0$ dan bahwa norma
(dan seminorma) hanya dapat mengambil nilai nonnegatif.
\end{exer}

Setiap ruang hasil kali dalam merupakan ruang linear bernorma. Dalam bahasa yang lebih tepat, hasil kali dalam pada suatu ruang vektor menginduksi
 \index{norma!diinduksi oleh hasil kali dalam}%
 \index{diinduksi!norma}
suatu norma pada ruang tersebut.

\begin{prop}\label{00015025} Misalkan $V$ suatu ruang hasil kali dalam. Pemetaan $x \mapsto \norm
x$ pada $V$ yang didefinisikan dalam~\ref{norm_notn} merupakan norma pada~$V$.
\end{prop}

Dan setiap ruang linear bernorma merupakan ruang metrik. Ingat kembali definisi berikut.

\begin{defn}\label{C019614} Misalkan $M$ suatu himpunan tak kosong. Sebuah fungsi $d \colon M \times M \sto \R$ adalah
 \index{metrik}%
\df{metrik} (atau
 \index{jarak!fungsi}%
 \index{d@$d(x,y)$ (jarak antara titik $x$ dan $y$)}%
\df{fungsi jarak}) pada $M$ jika untuk semua $x$, $y$, $z \in M$
 \begin{enumerate}
     \item $d(x,y) = d(y,x)$,
     \item $d(x,y) \le d(x,z) + d(z,y)$, \qquad (
 \index{ketaksamaan segitiga}%
 \index{ketaksamaan!segitiga}%
\emph{ketaksamaan segitiga})
     \item $d(x,x) = 0$, dan
     \item $d(x,y) = 0$ hanya jika $x = y$.
 \end{enumerate}
Jika $d$ merupakan metrik pada himpunan $M$, kita mengatakan bahwa pasangan $(M,d)$ merupakan
 \index{metrik!ruang}%
 \index{ruang!metrik}%
\df{ruang metrik}. Notasi ``$d(x,y)$'' dibaca ``jarak antara $x$
dan~$y$''. (Seperti biasa, kita mengganti perumusan yang benar ``Misalkan $(M,d)$ suatu ruang metrik \dots''
dengan ungkapan ``Misalkan $M$ suatu ruang metrik \dots''.)

Jika $d \colon M \times M \sto \R$ memenuhi syarat (1)--(3), tetapi tidak memenuhi (4), maka $d$ merupakan
 \index{pseudometrik}%
\df{pseudometrik} pada~$M$.
\end{defn}

Seperti disebutkan di atas, setiap ruang linear bernorma merupakan ruang metrik. Secara lebih tepat, norma pada suatu ruang vektor
menginduksi suatu
 \index{metrik!diinduksi oleh norma}%
metrik $d$, yang didefinisikan oleh $d(x,y) = \norm{x - y}$. Dengan kata lain, jarak antara
dua vektor adalah panjang selisih keduanya.
  \[ \xy
       \Atriangle/<-`<-`>/[``;x`x-y`y]
  \endxy \]
Jika tidak ada metrik lain yang ditentukan, kita selalu memandang ruang linear bernorma sebagai ruang metrik
 \index{diinduksi!metrik}%
di bawah metrik terinduksi ini. Metrik tersebut disebut \emph{metrik yang diinduksi oleh norma}.

\begin{exam}\label{0001503} Misalkan $V$ suatu ruang linear bernorma. Definisikan
$d\colon V \times V \sto \R$ dengan $d(x,y) = \norm{x - y}$. Maka $d$ merupakan metrik pada~$V$.
Jika $V$ hanya merupakan ruang berseminorma, maka $d$ merupakan pseudometrik.
\end{exam}

\begin{defn} Vektor $x$ dan $y$ dalam ruang hasil kali dalam $V$ disebut
 \index{ortogonal}%
\df{ortogonal} (atau
 \index{tegak lurus}%
\df{tegak lurus}) jika $\langle x,y \rangle = 0$. Dalam hal ini kita menulis
 \index{<binrelperp@$x \perp y$ (vektor ortogonal)}%
$x \perp y$. Subhimpunan $A$ dan $B$ dari $V$ disebut \df{ortogonal} jika $a \perp b$ untuk setiap $a
\in A$ dan $b \in B$. Dalam hal ini kita menulis $A \perp B$.
\end{defn}

\begin{defn} Jika $M$ dan $N$ merupakan subruang dari ruang hasil kali dalam $V$, kita menggunakan notasi
$V = M \oplus N$ untuk menyatakan bukan hanya bahwa $V$ merupakan jumlah langsung (ruang vektor) dari $M$
dan $N$, melainkan juga bahwa $M$ dan $N$ ortogonal. Dengan demikian, kita mengatakan bahwa $V$ merupakan
 \index{internal!jumlah langsung ortogonal}%
 \index{langsung!jumlah!ortogonal}%
 \index{ortogonal!jumlah langsung!ruang hasil kali dalam}%
\df{jumlah langsung ortogonal} (\df{internal}) dari $M$ dan~$N$.
\end{defn}

\begin{prop} Misalkan $a$ suatu vektor dalam ruang hasil kali dalam $V$. Maka $a \perp x$ untuk setiap
$x \in V$ jika dan hanya jika $a = 0$.
\end{prop}

\begin{prop} Dalam ruang hasil kali dalam, $x \perp y$ jika dan hanya jika $\norm{x + \alpha y} = \norm{x - \alpha y}$
untuk semua skalar $\alpha$.
\end{prop}

\begin{prop}[Teorema Pythagoras]\label{00015037} Jika $x \perp y$ dalam suatu ruang hasil
 \index{teorema Pythagoras}%
kali dalam, maka
   \[ \norm{x + y}^2 = \norm x^2 + \norm y^2\,. \]
\end{prop}

\begin{defn}\label{0001504} Misalkan $V$ dan $W$ merupakan ruang hasil kali dalam di atas medan skalar yang sama. Untuk $(v,w)$ dan
$(v',w')$ dalam $V \times W$ dan $\alpha \in \K$, definisikan
 \begin{align*}
      (v,w) + (v',w') &= (v + v', w + w') \\
\intertext{dan}
         \alpha(v,w) &= (\alpha v,\alpha w)\,.
 \end{align*}
Hasilnya adalah suatu ruang vektor, yaitu \emph{jumlah langsung (eksternal)} dari $V$ dan~$W$.
Untuk menjadikannya ruang hasil kali dalam, definisikan
   \[ \langle (v,w),(v',w') \rangle = \langle v,v' \rangle + \langle w,w' \rangle. \]
Definisi ini menjadikan jumlah langsung $V$ dan $W$ suatu ruang hasil kali dalam. Ruang tersebut merupakan
 \index{eksternal!jumlah langsung}%
 \index{langsung!jumlah!eksternal}%
 \index{langsung!jumlah!ruang hasil kali dalam}%
 \index{ortogonal!jumlah langsung!ruang hasil kali dalam}%
 \index{<binopdirectsum@$\oplus$ (jumlah langsung ortogonal)}%
\df{jumlah langsung} dari $V$ dan $W$ yang \df{ortogonal eksternal}, dan dinotasikan dengan $V \oplus W$.
\end{defn}

Perhatikan bahwa notasi $\oplus$ yang sama digunakan baik untuk jumlah langsung internal maupun eksternal,
serta baik untuk jumlah langsung ruang vektor (lihat definisi~\ref{000082}) maupun jumlah langsung ortogonal.
Jadi, ketika melihat simbol $V \oplus W$, penting untuk mengetahui kategori mana yang sedang kita gunakan:
ruang vektor atau ruang hasil kali dalam. Hal ini khususnya penting karena kata ``ortogonal'' lazim dihilangkan
sebagai pewatas ``jumlah langsung'', bahkan ketika sifat ortogonal memang dimaksudkan.

\begin{exam} Dalam $\R^2$, misalkan $M$ adalah sumbu-$x$ dan $L$ adalah garis dengan persamaan $y = x$.
Jika $\R^2$ dipandang sebagai ruang vektor (real), penulisan $\R^2 = M
\oplus L$ benar. Sebaliknya, jika $\R^2$ dipandang sebagai ruang hasil kali dalam (real), maka
$\R^2 \ne M \oplus L$ (karena $M$ dan $L$ tidak tegak lurus).
\end{exam}

\begin{prop} Misalkan $V$ suatu ruang hasil kali dalam. Hasil kali dalam pada $V$, yang dipandang sebagai pemetaan
dari $V \oplus V$ ke $\K$, bersifat kontinu. Demikian pula normanya, yang dipandang sebagai pemetaan dari $V$
ke $\R$.
\end{prop}

Mengenai pembuktian proposisi di atas, perhatikan bahwa pemetaan $(v,v') \mapsto
\norm v + \norm{v'}$, $(v,v') \mapsto \sqrt{\norm v^2 + \norm{v'}^2}$, dan $(v,v')
\mapsto \max\{\norm v, \norm{v'}\}$ semuanya merupakan norma pada $V \oplus V$. Norma manakah yang diinduksi
oleh hasil kali dalam pada $V \oplus V$? Mengapa pilihan norma tidak berpengaruh ketika kita membuktikan
bahwa hasil kali dalam tersebut kontinu?

\begin{prop}[Hukum jajaran genjang]\label{parallelogram_law} Jika $x$ dan $y$ merupakan vektor dalam suatu ruang hasil kali
 \index{hukum jajaran genjang}%
dalam, maka
   \[ \norm{x + y}^2 + \norm{x - y}^2 = 2\norm x^2 + 2\norm y^2\,. \]
\end{prop}

Walaupun setiap hasil kali dalam menginduksi suatu norma, tidak setiap norma berasal dari hasil kali dalam.

\begin{exam}\label{00015061} Tidak ada hasil kali dalam pada ruang $\fml C([0,1])$ yang menginduksi norma seragam.
\end{exam}

\begin{proof}[\emph{Petunjuk pembuktian}] Gunakan proposisi sebelumnya. \ns
\end{proof}

\begin{prop}[Identitas polarisasi]\label{0001507} Jika $x$ dan $y$ merupakan vektor
 \index{identitas polarisasi}%
dalam ruang hasil kali dalam kompleks, maka
 \[ \langle x,y \rangle = \tfrac14 (\norm{x + y}^2 - \norm{x - y}^2
                       + i\,\norm{x + iy}^2 - i\,\norm{x - iy}^2)\,. \]
\end{prop}

Apakah identitas yang benar untuk ruang hasil kali dalam \emph{real}?

\begin{notn}\label{0001511} Misalkan $V$ suatu ruang hasil kali dalam, $x \in V$, dan $A$,
$B \subseteq V$. Jika $x \perp a$ untuk setiap $a \in A$, kita menulis $x \perp A$; dan jika $a
\perp b$ untuk setiap $a \in A$ dan $b \in B$, kita menulis $A \perp B$.
Kita mendefinisikan $A^\perp$, yaitu
 \index{ortogonal!komplemen}%
 \index{komplemen!ortogonal}%
 \index{<unopur@$A^\perp$ (komplemen ortogonal)}%
\df{komplemen ortogonal} dari $A$, sebagai $\{x \in V\colon x \perp A\}$. Karena frasa ``komplemen ortogonal dari $A$''
agak panjang, terutama jika digunakan berulang kali, simbol ``$A^\perp$'' biasanya dilafalkan ``A perp''.
Kita menulis $A^{\perp\perp}$ untuk $\left(A^\perp\right)^\perp$.
\end{notn}

\begin{prop}\label{0001512} Jika $A$ merupakan subhimpunan dari ruang hasil kali dalam $V$, maka
$A^\perp$ merupakan subruang linear tertutup dari~$V$.
\end{prop}

\begin{prop}[Ortonormalisasi Gram--Schmidt] Jika $(v^k)$ merupakan barisan bebas linear dalam
 \index{ortonormalisasi Gram--Schmidt}%
 \index{ortonormalisasi}%
ruang hasil kali dalam~$V$, maka terdapat barisan ortonormal $(e^k)$ dalam~$V$ sedemikian sehingga
$\spn\{v^1, \dots, v^n\} = \spn\{e^1, \dots, e^n\}$ untuk setiap $n \in \N$.
\end{prop}

\begin{cor}\label{0001514} Jika $M$ merupakan subruang dari ruang hasil kali dalam berdimensi
hingga $V$, maka $V = M \oplus M^\perp$.
\end{cor}

Akan berguna untuk mengatakan bahwa suatu barisan \emph{pada akhirnya} memiliki sifat tertentu atau
bahwa barisan itu \emph{sering} memiliki sifat tersebut.

\begin{defn}\label{C013121} Misalkan $(x_n)$ suatu barisan elemen dari himpunan $S$ dan $P$ suatu
sifat yang mungkin dimiliki anggota $S$. Kita mengatakan bahwa barisan $(x_n)$
 \index{pada akhirnya}%
\df{pada akhirnya} memiliki sifat $P$ jika terdapat $n_0 \in \N$ sedemikian sehingga $x_n$ memiliki
sifat $P$ untuk setiap $n \ge n_0$. (Cara lain untuk menyatakan hal yang sama: $x_n$ memiliki
sifat $P$ untuk semua kecuali sejumlah hingga~$n$.)
\end{defn}

\begin{exam}\label{exam13122} Misalkan $l_c$ menyatakan
 \index{l@$l_c$!barisan yang pada akhirnya nol}%
 \index{l@$l_c$!sebagai ruang vektor}%
ruang vektor yang terdiri atas semua barisan bilangan kompleks yang pada akhirnya nol;
yaitu, himpunan semua barisan yang hanya memiliki sejumlah hingga entri tak nol. Barisan tersebut juga disebut barisan dengan
 \index{hingga!tumpuan}%
 \index{tumpuan!hingga}%
\emph{bertumpuan hingga}. Operasi ruang vektornya didefinisikan titik demi titik. Kita menjadikan ruang $l_c$ suatu
 \index{l@$l_c$!sebagai ruang hasil kali dalam}%
 \index{hasil kali dalam!ruang!$l_c$ sebagai}%
ruang hasil kali dalam dengan
mendefinisikan $\langle a,b \rangle = \langle\, (a_n),(b_n)\, \rangle := \sum_{k=1}^\infty a_k \conj{b_k}$.
\end{exam}

\begin{defn}\label{C013131} Misalkan $(x_n)$ suatu barisan elemen dari himpunan $S$ dan $P$ suatu
sifat yang mungkin dimiliki anggota $S$. Kita mengatakan bahwa barisan $(x_n)$
 \index{sering}%
\df{sering} memiliki sifat $P$ jika untuk setiap $k \in \N$ terdapat $n \ge k$ sedemikian sehingga
$x_n$ memiliki sifat~$P$. (Perumusan yang ekuivalen: $x_n$ memiliki sifat $P$ untuk tak hingga banyak
nilai~$n$.)
\end{defn}

\begin{exam}\label{exam13132} Misalkan $V = l_2$ merupakan ruang hasil kali dalam semua barisan bilangan kompleks
yang terjumlah secara kuadrat (lihat contoh~\ref{000150013}) dan $M = l_c$ (lihat contoh~\ref{exam13122}). Maka kesimpulan~\ref{0001514} gagal.
\end{exam}

\begin{defn}\label{def_alg_dual} Suatu
 \index{linear!fungsional}%
 \index{fungsional!linear}%
\df{fungsional linear} pada ruang vektor $V$ adalah pemetaan linear dari $V$ ke medan
skalarnya. Himpunan semua fungsional linear pada $V$ merupakan
 \index{dual!aljabar}%
 \index{aljabar!ruang dual}%
 \index{ruang!dual aljabar}%
\df{ruang dual} (\df{aljabar}) dari~$V$. Kita akan menggunakan
 \index{<unopur@$V^\#$ (ruang dual aljabar)}%
 \index{V@$V^\#$ (ruang dual aljabar)}%
notasi $V^{\#}$ untuk ruang dual aljabar.
\end{defn}

\begin{thm}[Teorema Riesz--Fr\'echet]\label{0001601} Jika $f \in V^{\#}$ dengan $V$ suatu
 \index{teorema Riesz--Fr\'echet!versi berdimensi hingga}%
ruang hasil kali dalam berdimensi hingga, maka terdapat vektor unik $a$ dalam $V$ sedemikian
sehingga
    \[ f(x) = \langle x,a \rangle \]
untuk semua $x$ dalam $V$.
\end{thm}

Sebentar lagi kita akan membuktikan (dalam Teorema~\futurexref{4.5.2}{000342}) bahwa setiap fungsional linear kontinu pada ruang hasil
kali dalam \emph{sebarang} memiliki representasi di atas. Versi berdimensi hingga yang dinyatakan di sini merupakan kasus khusus, karena
setiap pemetaan linear pada ruang hasil kali dalam berdimensi hingga bersifat kontinu (lihat Proposisi~\futurexref{3.3.9}{00015033}).

\begin{defn}\label{0001602} Misalkan $T\colon V \sto W$ suatu transformasi linear antara ruang
hasil kali dalam kompleks. Jika terdapat fungsi $T^*\colon W \sto V$ yang memenuhi
   \[ \langle T\vc v,\vc w \rangle = \langle \vc v,T^*\vc w \rangle \]
untuk semua $\vc v \in V$ dan $\vc w \in W$, maka $T^*$ merupakan
 \index{adjoin!operator pada ruang hasil kali dalam}%
 \index{<unopurstar@$T^*$ (adjoin pada ruang hasil kali dalam)}%
 \index{T@$T^*$ (adjoin pada ruang hasil kali dalam)}%
\df{adjoin} (atau \df{transpos konjugat}, atau \df{konjugat Hermitian}) dari~$T$.
\end{defn}

\begin{prop}\label{00016021} Jika $T\colon V \sto W$ merupakan pemetaan linear antara ruang
hasil kali dalam berdimensi hingga, maka $T^*$ ada.
\end{prop}

\begin{proof}[\emph{Petunjuk pembuktian}] Fungsional $\phi\colon V \times W \sto \C \colon
(v,w) \mapsto \langle Tv,w \rangle$ bersifat seskuilinear. Tetapkan $w \in W$ dan definisikan $\phi_w
\colon V \sto \C \colon v \mapsto \phi(v,w)$. Maka $\phi_w \in V^{\#}$. Gunakan
\emph{teorema Riesz--Fr\'echet}~(\ref{0001601}). \ns
\end{proof}

\begin{prop} Jika $T\colon V \sto W$ merupakan pemetaan linear antara ruang hasil kali dalam
berdimensi hingga, maka fungsi $T^*$ yang didefinisikan di atas bersifat linear dan $T^{**} = T$.
\end{prop}

\begin{thm}[Teorema dasar aljabar linear]\label{00016025} Jika
 \index{dasar!teorema aljabar linear}%
$T\colon V \sto W$ merupakan pemetaan linear antara ruang hasil kali dalam berdimensi hingga, maka
   \[ \ker T^* = (\ran T)^\perp \quad \text{ dan } \quad \ran T^* = (\ker T)^\perp\,. \]
\end{thm}

\begin{defn}\label{000161} Suatu operator $U$ pada ruang hasil kali dalam disebut
 \index{uniter!operator}%
 \index{operator!uniter}%
\df{uniter} jika $UU^* = U^*U = I$, yaitu jika $U^* = U^{-1}$.
\end{defn}

\begin{defn}\label{0001612} Dua operator $R$ dan $T$ pada ruang hasil kali dalam $V$ disebut
 \index{uniter!ekuivalensi}%
 \index{ekuivalen!secara uniter}%
\df{ekuivalen secara uniter} jika terdapat operator uniter $U$ pada $V$ sedemikian sehingga $R =
U^*TU$.
\end{defn}

\begin{prop} Jika $V$ merupakan ruang hasil kali dalam, maka ekuivalensi uniter memang merupakan
relasi ekuivalensi pada~$\ofml L(V)$.
\end{prop}

\begin{defn} Suatu operator $T$ pada ruang hasil kali dalam berdimensi hingga $V$ disebut
 \index{uniter!diagonalisasi}%
 \index{dapat didiagonalkan!secara uniter}%
\df{dapat didiagonalkan secara uniter} jika terdapat basis ortonormal untuk $V$ yang terhadapnya
$T$ bersifat diagonal. Secara ekuivalen, suatu operator pada ruang hasil kali dalam berdimensi hingga
\emph{dengan basis} dapat didiagonalkan secara uniter jika operator itu ekuivalen secara uniter dengan operator
diagonal.
\end{defn}

\begin{defn} Suatu operator $T$ pada ruang hasil kali dalam disebut
 \index{swaadjoin!operator}%
 \index{operator!swaadjoin}%
\df{swaadjoin} (atau
 \index{Hermitian!operator}%
 \index{operator!Hermitian}%
\df{Hermitian}) jika $T^* = T$.
\end{defn}

\begin{defn}\label{000164} Suatu proyeksi $P$ dalam ruang hasil kali dalam disebut
 \index{proyeksi!ortogonal}%
 \index{ortogonal!proyeksi!dalam ruang hasil kali dalam}%
 \index{operator!proyeksi ortogonal}%
\df{proyeksi ortogonal} jika proyeksi itu swaadjoin. Jika $M$ merupakan jangkauan suatu proyeksi
ortogonal, kita akan menggunakan notasi
 \index{PM@$\sbsb PM$ (proyeksi ortogonal ke~$M$)}%
$\sbsb PM$ untuk proyeksi tersebut, bukan notasi yang lebih rumit~$\sbsb E{M^\perp M}$.
\end{defn}

\begin{cau} Proyeksi pada ruang vektor atau ruang linear bernorma bersifat linear dan idempoten,
sedangkan proyeksi ortogonal pada ruang hasil kali dalam bersifat linear, idempoten, dan
swaadjoin. Situasi yang sebenarnya langsung ini agak dirumitkan oleh kecenderungan umum
untuk menyebut proyeksi ortogonal hanya sebagai ``proyeksi''. Bahkan, kelak dalam catatan
ini kita akan menggunakan konvensi tersebut. Dalam ruang hasil kali dalam,
$\oplus$ biasanya menyatakan jumlah langsung ortogonal dan ``proyeksi'' biasanya berarti
``proyeksi ortogonal''. Dalam banyak buku pengantar aljabar linear, seluruh pembahasan
berlangsung dalam $\R^n$, sehingga kadang sulit menentukan apakah pada suatu halaman
tertentu penulis memperlakukan $\R^n$ sebagai ruang vektor atau sebagai ruang hasil kali
dalam.
\end{cau}

\begin{prop} Jika $P$ merupakan proyeksi ortogonal pada ruang hasil kali dalam $V$, maka kita memperoleh
dekomposisi jumlah langsung ortogonal $V = \ker P \oplus \ran P$.
\end{prop}

\begin{defn} Jika $I_V = P_1 + P_2 + \dots + P_n$ merupakan resolusi identitas dalam suatu
ruang hasil kali dalam~$V$ dan setiap $P_k$ merupakan proyeksi ortogonal, kita mengatakan bahwa $I_V =
P_1 + P_2 + \dots + P_n$ merupakan
 \index{resolusi identitas!ortogonal}%
 \index{ortogonal!resolusi identitas}%
\df{resolusi ortogonal dari identitas}.
\end{defn}

\begin{prop} Jika $I_V = P_1 + P_2 + \dots + P_n$ merupakan resolusi ortogonal dari identitas
pada ruang hasil kali dalam $V$, maka $V = \bigoplus_{k=1}^n \ran P_k$.
\end{prop}

\begin{defn} Suatu operator $N$ pada ruang hasil kali dalam disebut
 \index{normal!operator}%
 \index{operator!normal}%
\df{normal} jika $NN^* = N^*N$.
\end{defn}

Dua pencapaian besar aljabar linear adalah \emph{teorema spektral} untuk operator pada
ruang hasil kali dalam (kompleks) berdimensi hingga (lihat \ref{00017}), yang memberikan
syarat perlu dan cukup yang dapat dinyatakan secara sederhana agar suatu operator dapat
didiagonalkan secara uniter, dan teorema~\ref{00018}, yang memberikan klasifikasi lengkap
operator-operator tersebut.

\begin{thm}[Teorema Spektral untuk Ruang Hasil Kali Dalam Kompleks Berdimensi Hingga]\label{00017}
Misalkan $T$ suatu operator pada ruang hasil kali dalam berdimensi hingga~$V$ dengan nilai-nilai eigen (berbeda)
 \index{spektral!teorema!operator normal pada ruang hasil kali dalam}%
$\lambda_1, \dots, \lambda_n$. Maka $T$ dapat didiagonalkan secara uniter jika dan
hanya jika $T$ normal. Jika $T$ normal, maka terdapat resolusi ortogonal dari identitas
$I_V = P_1 + \dots + P_n$, dengan jangkauan proyeksi ortogonal $P_k$ untuk setiap $k$
merupakan ruang eigen yang berkaitan dengan~$\lambda_k$, dan selanjutnya
     \[ T = \lambda_1P_1 +\dots + \lambda_nP_n\,. \]
\end{thm}

\begin{proof} Lihat \cite{Roman:2005}, halaman 227. \ns
\end{proof}

\begin{thm}\label{00018} Dua operator normal pada ruang hasil kali dalam berdimensi hingga
 \index{normal!operator!syarat untuk ekuivalensi uniter}%
ekuivalen secara uniter jika dan hanya jika keduanya memiliki nilai-nilai eigen yang sama,
masing-masing dengan multiplisitas yang sama; yaitu, jika dan hanya jika keduanya memiliki
polinom karakteristik yang sama.
\end{thm}

\begin{proof} Lihat \cite{HoffmanK:1971}, halaman 357. \ns
\end{proof}

Sebagian besar uraian selanjutnya dikhususkan untuk upaya yang sungguh tak sepele dalam mencari
generalisasi yang tepat dari kedua hasil di atas ke ranah berdimensi tak hingga, serta bentang
matematika menakjubkan yang mulai terlihat di sepanjang perjalanan tersebut.

\begin{exer} Misalkan $N = \dfrac13\begin{bmatrix}
                                4+2i & 1-i  & 1-i  \\
                                1-i  & 4+2i & 1-i  \\
                                1-i  & 1-i  & 4+2i \end{bmatrix}$.

\vskip 3 pt
 \begin{enumerate}
  \item Tunjukkan bahwa matriks $N$ bersifat normal.

\vskip 3 pt

  \item Dengan demikian, menurut \emph{teorema spektral}, $N$
dapat dituliskan sebagai kombinasi linear proyeksi-proyeksi ortogonal.
Tunjukkan secara eksplisit cara melakukannya (dan kerjakan perhitungannya).
 \end{enumerate}
\end{exer}

















\endinput
